\section{Bounded loops and quantifier-free induction}

\subsection{The terminal-extraction pattern}

Let \(I(t,\vec x)\) be QF, and suppose iPRA proves
\[
 I(0,\vec x),
 \qquad
 I(t,\vec x)\to I(t+1,\vec x).
\]
Quantifier-free induction yields \(I(T(\vec x),\vec x)\).  If a separate object
proof gives
\[
 I(T(\vec x),\vec x)\to C(\vec x),
\]
then \(C\) may contain arbitrary quantifiers and connectives.

\begin{proposition}[QF loop, arbitrary finish]\label{prop:qf-loop}
A bounded computation needs only a QF running invariant.  Its terminal theorem
need not be QF.
\end{proposition}

\begin{proof}
Apply \leanname{PRA.qfInd}, eliminate the universal quantifier at the explicit
fuel term, and use implication elimination with the terminal extractor.  This
is the theorem \leanname{loopProduces} in the supplied structural core.
\end{proof}

\subsection{Fixed precision, quantified last}

An analytic packet often wants to prove
\[
  \forall p\;\exists c\;\Valid(p,c).
\]
Induction over the whole displayed formula is unavailable because it contains
quantifiers.  The correct order is:

\begin{enumerate}
\item treat \(p\) as a free parameter;
\item define a QF invariant for the finite computation at that precision;
\item prove the terminal checker equation;
\item introduce the certificate existential or selected term;
\item universally quantify \(p\) after the induction.
\end{enumerate}

The same pattern handles a theorem with several precision levels: pack them
into the finite query first.

\subsection{Changing parameters}

An iterative algorithm may change not only its state but also its current
precision, denominator bound, active chart, or remaining resource.  A naive
invariant with a fixed external parameter is then false.  The remedy is to
include every changing value in the recursion state.

The repository's pairing and simultaneous-recursion modules give two standard
encodings:

\begin{itemize}
\item pack a tuple into one natural state and update all fields;
\item use simultaneous primitive recursion for a small fixed number of
sequences.
\end{itemize}

The QF invariant is stated on the packed current state, time, and immutable
query parameters.  The checked object theorem
\leanname{primitiveRecursionInvariant} was designed for exactly this use.

\subsection{Trace certificates versus computed traces}

There are two ways to handle a bounded loop.

\paragraph{Computed trace.}
A fixed PR transition computes the final state.  Prove an invariant by
primitive recursion.  This is compact and efficient.

\paragraph{Certified trace.}
The entire sequence of states or branch bits is a single-code array.  A checker
verifies the initial state, every transition, and the terminal condition.
Existence may be established by a producer or by packing the computed trace.
This version is more flexible when a host tool supplies a trace or when the
transition is relational.

The two forms are interderivable when the transition is deterministic.  A
useful engineering policy is to prove the certified-trace checker first, then
show that the efficient computed trace is accepted.

\subsection{Lexicographic descent}

Nested or lexicographic descent can be reduced to a single natural fuel.  Given
bounded components \(a_i<M_i(q)\), use their mixed-radix rank
\[
  \rho(\vec a)=a_0+M_0a_1+M_0M_1a_2+\cdots.
\]
A lexicographic decrease becomes a strict decrease of \(\rho\) after choosing
the digit order appropriately.  Primitive recursion for at most
\(\prod_iM_i(q)\) stages then proves termination.

\begin{warning}
A qualitative statement that a well-founded measure exists is insufficient.
The rank and its numerical bound must be explicit PR terms, and the decrease
must be an object-iPRA inequality.
\end{warning}
